\section{De la H\&E de bajo costo a representaciones moleculares ricas}

La sección anterior usó la mask y el contexto espacial dentro de un mismo tipo de datos genéticos. Esta sección vuelve a la translation: ¿es posible predecir mediciones de transcriptómica espacial o de proteínas, costosas y escasas, a partir de imágenes de H\&E, baratas y comunes en la práctica clínica? Si funciona, el modelo puede convertir los datos históricos de patología en representaciones más ricas de los pacientes; si falla, puede producir pseudomapas moleculares sumamente engañosos.

\subsection{La tarea de entrenamiento de la imputación multimodal}

La imputation (imputación) consiste en estimar variables faltantes a partir de las variables observadas. No es copiar el valor promedio, sino aprender una distribución condicional, y debe hacer explícita la incertidumbre de la predicción. Aquí la modalidad faltante es un molecular assay costoso. A continuación se examina, desde tres ángulos (el eje de costo, la arquitectura condicional y la mask de la modalidad completa), si el modelo realmente reproduce las condiciones de información del despliegue clínico, donde solo se dispone de H\&E.

\lecturefigure{slide-213-multimodal-imputation.jpg}{Inferir modalidades moleculares costosas y escasas a partir de la H\&E, barata y común}{00:43:48}

El eje de costo de la figura muestra el valor de la tarea: la H\&E es fácil de obtener y las mediciones moleculares espaciales son costosas. Si el modelo solo se entrena con datos pareados pero se despliega en otros hospitales, con otros protocolos de tinción u otros tipos de tumor, la translation puede colapsar. Por eso el valor proviene de la generalización entre dominios y de la calibración, no solo de reconstruir muestras pareadas.

\lecturefigure{slide-215-he-conditioned-model.jpg}{Modelo masked gene-count condicionado por un codificador de H\&E}{00:44:27}

La arquitectura codifica cada patch de H\&E como condición y la combina con la red troncal de tokens de genes. El H\&E-specific encoder se encarga de las características visuales y la red troncal, de la count prediction. Esta división permite que el mismo modelo genético funcione con o sin imagen, pero también puede llevar al codificador visual a aprender atajos basados en el lote de tinción.

\lecturefigure{slide-218-cross-modal-translation.jpg}{Aplicar mask a toda la entrada genética obliga al modelo a predecir las lecturas moleculares solo a partir de la H\&E}{00:45:00}

Cuando se aplica mask a todos los gene tokens, la salida depende por completo de la condición de H\&E. En ese caso la tarea pasa de disambiguation a una translation más pura. Si se conservan algunas lecturas genéticas reales, el modelo se desplaza entre translation y disambiguation: la imagen aporta una distribución a priori morfológica y los marker medidos corrigen la incertidumbre.

\begin{importantbox}{El objetivo de entrenamiento debe simular lo que estará disponible en el despliegue}
Si el despliegue clínico solo cuenta con H\&E, la evaluación debe ocultar toda la entrada molecular; si el despliegue admite algunos marker, debe medirse la curva de desempeño con distintos presupuestos de marker. Ver la modalidad molecular completa durante el entrenamiento produce un techo imposible de alcanzar en el despliegue.
\end{importantbox}

\subsection{Mapas de predicción y estructura del error}

El resultado final no puede evaluarse con un solo mapa de calor de colores parecidos. Hay que comparar la medición real, la baseline de interpolación y la predicción del modelo, y desglosar el error por gen, paciente, región y tipo celular. Además, conviene preguntar si los estados moleculares que no son identificables por la morfología corresponden a intervalos de incertidumbre más amplios, y si el modelo se mantiene calibrado entre hospitales y entre lotes de tinción.

\lecturefigure{slide-224-gene-expression-prediction.jpg}{Predicción, a partir de H\&E, de tipos celulares y de la expresión espacial de genes como KRT19}{00:46:28}

A la izquierda está la H\&E; en el centro, los tipos celulares o el overlay; a la derecha, la expresión predicha. Primero se compara si las estructuras a gran escala coinciden y después se revisan los bordes, las regiones poco frecuentes y los picos locales. Cuando la morfología y la expresión están correlacionadas, el modelo puede desempeñarse muy bien, pero los estados moleculares distintos con la misma morfología pueden seguir siendo imposibles de determinar a partir de la H\&E.

\lecturefigure{slide-225-imputation-accuracy.jpg}{La dificultad de imputación varía notablemente entre genes}{00:47:16}

Las tres filas comparan la detección real, la interpolación y la predicción del modelo. Algunos genes tienen una correspondencia morfológica clara y se predicen bien; otros casi no pueden determinarse de forma única a partir de la H\&E. Las métricas promedio ocultan esta diferencia, por lo que se necesita gene-wise calibration: el modelo debe saber en qué objetivos carece de información identificable.

Si el valor real es $y$, la media de la distribución predicha es $\hat y$ y la desviación estándar es $\hat\sigma$, la calibración no solo exige un error pequeño, sino también
\[
\Pr\left(|y-\hat y|\leq 1.96\hat\sigma\right)\approx 0.95.
\]
La cobertura de los intervalos debe reportarse por separado para held-out patient, para distintos centros y para distintos genes.

\begin{warningbox}{La predictibilidad a partir de la morfología tiene un límite}
Si dos regiones de tejido son casi idénticas en la H\&E pero tienen estados de RNA distintos, ningún modelo que solo vea la H\&E podrá distinguirlas de forma confiable. En ese caso, la salida correcta es una incertidumbre alta, no un mapa de calor suavizado generado con el patrón promedio del conjunto de entrenamiento.
\end{warningbox}

\lecturefigure{slide-234-representation-pipeline.jpg}{De la H\&E a representaciones ricas imputadas por el modelo, y de ahí a hipótesis terapéuticas}{00:50:01}

Esta figura de síntesis deja claro el valor de ingeniería de la translation: primero se construye una representación de alta dimensión del paciente a partir de la H\&E de bajo costo y después se usa para búsqueda, estratificación o hipótesis farmacológicas. Cualquier paso intermedio puede amplificar los sesgos. La imputación del modelo no es un sustituto sin pérdida de la medición real; se parece más a un assay virtual escalable y con incertidumbre.

\teachervoice{00:43:46--00:48:40: el docente enfatiza que la calidad de la predicción no es uniforme entre genes y muestra la detección real, la interpolación y la salida del modelo. La conclusión de la clase no es «recuperar todo a partir de la H\&E», sino identificar qué información molecular puede quedar restringida de forma confiable por la morfología.}

\begin{knowledgebox}{Primera aparición: Calibration}
La calibration (calibración) exige que la confianza coincida con la tasa real de aciertos. Si el modelo sigue dando intervalos estrechos para genes no identificables, aunque el error promedio sea aceptable, inducirá a error al priorizar experimentos. En AI for Science, «no saber que no se sabe» suele ser más peligroso que un error grande ocasional.
\end{knowledgebox}

\subsection{Resumen de la sección}

La imputation entre modalidades convierte la ventaja de escala de la H\&E en hipótesis moleculares, pero su límite lo fija la información identificable. Un sistema confiable debe comparar, con particiones a nivel de paciente, la medición real, una baseline simple y la predicción del modelo; reportar el error y la calibración por gen, y tratar los resultados imputados como un assay virtual, no como un experimento real.
